\label{sec:design}

\ssrov{} uses one holder-signed policy containing a conventional ROA entry and nonoverlapping prefix markers. Each marker applies only to proper subprefixes no longer than the policy's maximum route-prefix length.

\subsection{\ssrov{} Policy}

A policy is
\begin{equation}
D=(v,p_0,m,a_h,\mathcal{R}),
\end{equation}
where $v$ is the profile version, $p_0$ the base prefix, $m$ the maximum route-prefix length, and $a_h$ the known aggregate origin. The policy carries the conventional ROA entry $(p_0,a_h,m)$; $\mathcal{R}$ contains the prefix markers without an origin list.

The current semantics use profile version $v=2$; the address family follows $p_0$, with width $w\in\{32,128\}$ and eligible ASN set $\mathcal{A}$. A valid policy has $|p_0|\leq m\leq w$; every $q\in\mathcal{R}$ is canonical, shares the address family of $p_0$, satisfies $q\subseteq p_0$ and $|q|<m$; and prefixes in $\mathcal{R}$ are disjoint. The tested profile excludes AS~0~\cite{rfc7607}, AS\_TRANS and out-of-range values~\cite{rfc6793}, private~\cite{rfc6996}, documentation~\cite{rfc5398}, and reserved~\cite{iana-as-numbers} values from $\mathcal{A}$; unsupported profiles and mixed address families are invalid.

A marker $q$ covers only proper subprefixes $p\subset q$ with $|p|\leq m$. Marker disjointness makes membership unique; ordinary ROAs may still cover these prefixes.

\subsection{Classification with Provenance}

Let $\rho(p,o)$ be the conventional ROV state of route $(p,o)$ under the validated ROAs, including $(p_0,a_h,m)$. Define the marker-eligibility predicate
\begin{equation}
\mathsf{Scope}_D(p,o)\iff
\begin{aligned}[t]
&|p|\leq m\land o\in\mathcal{A}\setminus\{a_h\}\land{}\\[-2pt]
&\exists q\in\mathcal{R}:p\subset q.
\end{aligned}
\end{equation}
For a valid policy $D$, the complete four-state classifier is
\begingroup
\small
\begin{equation}
\mathsf{State}_D(p,o)=
\begin{cases}
\explicit, & \rho(p,o)=\valid,\\
\openscope, &
\begin{aligned}
&\rho(p,o)=\invalid\\[-2pt]
&{}\land\mathsf{Scope}_D(p,o),
\end{aligned}\\
\invalid, &
\begin{aligned}
&\rho(p,o)=\invalid\\[-2pt]
&{}\land\neg\mathsf{Scope}_D(p,o),
\end{aligned}\\
\notfound, & \rho(p,o)=\notfound.
\end{cases}
\end{equation}
\endgroup
The cases are mutually exclusive and exhaustive, and an exact ROA takes precedence: a matching ROA yields \explicit{}, and an \invalid{} route becomes \openscope{} exactly when $\mathsf{Scope}_D$ holds. The output preserves provenance, because \explicit{} records matching ROA evidence while \openscope{} records marker, length, and ASN eligibility without origin authorization; local policy may accept, deprefer, or reject \openscope{} routes.

\subsection{Properties}

For the aggregate--subprefix dilemma, strict expressiveness requires two properties. With $\mathcal{R}=\varnothing$, \ssrov{} preserves conventional ROV. With markers present, it can return \openscope{} when $\mathsf{Scope}_D$ holds and no ROA matches the downstream origin.

Setting $\mathcal{R}=\varnothing$ gives the preservation condition: the classifier maps conventional \valid{}, \invalid{}, and \notfound{} to \explicit{}, \invalid{}, and \notfound{}, so a conventional entry $(q,a,\ell)$ maps to
\begin{equation}
D=(v,q,\ell,a,\varnothing).
\end{equation}
Discarding $\mathcal{R}$ from any policy leaves its conventional entry $(p_0,a_h,m)$ unchanged, and separately published exact ROAs need no projection to keep providing explicit origin matches.

With $\mathcal{R}$ nonempty, \openscope{} applies only when $p\subset q$ and $|p|\leq m$ for some $q\in\mathcal{R}$ and $o\in\mathcal{A}\setminus\{a_h\}$. Disjointness makes membership unambiguous, and an exact ROA below the same marker still yields \explicit{} because conventional matching comes first.

The following proposition establishes this separation.

\begin{proposition}[Aggregate--subprefix state separation]
\label{prop:scoped-state-separation}
Let
\[
D_{\mathsf{scope}}=(v,p_0,m,a_h,\{p_0\})
\]
be a valid policy. For every proper subprefix $p\subset p_0$ with $|p|\leq m$ and every $o\in\mathcal{A}\setminus\{a_h\}$ for which no ROA matches $(p,o)$,
\[
\begin{aligned}
\mathsf{State}_{D_{\mathsf{scope}}}(p_0,a_h)&=\explicit,\\
\mathsf{State}_{D_{\mathsf{scope}}}(p,o)&=\openscope.
\end{aligned}
\]
The construction is independent of which eligible downstream origin is selected, and the policy does not list that origin.
\end{proposition}

\begin{proof}
The conventional entry $(p_0,a_h,m)$ matches the aggregate route, so its state is \explicit{}. For $(p,o)$, the same entry provides coverage, while the premise excludes a matching ROA. Proposition~\ref{prop:aggregate-subprefix-dilemma} therefore gives the conventional state \invalid{}. Because $p\subset p_0$, $|p|\leq m$, and $o$ is eligible, the marker at $p_0$ satisfies $\mathsf{Scope}_{D_{\mathsf{scope}}}(p,o)$ and the classifier returns \openscope{}.
\end{proof}

The empty marker mapping and Proposition~\ref{prop:scoped-state-separation} satisfy both conditions. Conventional ROV has no state corresponding to \openscope{}, and Proposition~\ref{prop:aggregate-subprefix-dilemma} makes an unmatched proper subprefix with a different origin \invalid{}. Guaranteeing \valid{} for whichever of two indistinguishable origins is selected requires matching both, which Proposition~\ref{prop:origin-choice-visibility} shows also matches the unselected origin. \ssrov{} instead records markers and a length bound without enumerating either origin, so it is strictly more expressive for this dilemma in the stated sense.

\subsection{Security Scope and Threat Model}
\label{sec:security-scope}

{\revstart We consider an adversary that announces routes from AS numbers it controls and knows the policy and its markers, but cannot forge the holder's signature, alter validated RPKI data, or modify the policy. Choosing a prefix $p$ below a marker $q\in\mathcal{R}$ and an eligible origin $o\in\mathcal{A}\setminus\{a_h\}$ with no matching ROA yields \openscope{} instead of \invalid{}, and the receiving network applies its local action. Outside the markers, above $m$, or with an ineligible ASN the route remains \invalid{}, so the adversary gains nothing there.}

{\revstart Three properties bound the exposure: the holder chooses the marked region, so $q\subset p_0$ keeps it inside the aggregate and disjoint markers narrow it; the length bound caps any admitted route; and the eligible ASN set excludes AS~0, AS\_TRANS, and private, documentation, and reserved values. The holder removes the exposure by issuing a matching ROA, which takes precedence, and by shrinking or revoking $\mathcal{R}$; the ROA alone leaves other origins eligible below a retained marker.}

{\revstart The maximum route-prefix length is shared with the conventional entry by construction, so \openscope{} never admits a route longer than the length already authorized for $a_h$: the policy relaxes exactly one field, the origin constraint, and leaves the length envelope unchanged. A separate length bound for the marker branch would admit lengths outside that envelope; we do not evaluate that variant and leave it as an open design choice.}

Section~\ref{sec:implementation} implements this classifier in a reference policy engine and connects validated policies to router enforcement.
